\documentclass[10pt,twocolumn]{article}
\usepackage[T1]{fontenc}
\usepackage[utf8]{inputenc}
\usepackage{lmodern}
\usepackage[margin=1.8cm,top=2.4cm,bottom=2cm,columnsep=0.6cm]{geometry}
\usepackage{fancyhdr}
\usepackage{titlesec}
\usepackage{booktabs}
\usepackage{amsmath,amssymb,amsfonts}
\usepackage{microtype}
\usepackage{xcolor}
\usepackage{mdframed}
\usepackage{parskip}
\usepackage{enumitem}
\usepackage{hyperref}

\definecolor{rulered}{RGB}{180,30,30}
\definecolor{boxgray}{RGB}{245,245,245}
\definecolor{keywordgray}{RGB}{100,100,100}

\pagestyle{fancy}
\fancyhf{}
\fancyhead[L]{\textsc{\small Claude Mag}}
\fancyhead[C]{\small\textit{Machine Learning Research Digest}}
\fancyhead[R]{\small 2026-10-05}
\fancyfoot[C]{\small\thepage}
\renewcommand{\headrulewidth}{0.8pt}

\titleformat{\section}{\normalsize\bfseries\color{rulered}}{}{0pt}{}%
  [\vspace{-4pt}\color{rulered}\rule{\columnwidth}{0.4pt}]
\titlespacing{\section}{0pt}{8pt}{4pt}

\hypersetup{colorlinks=true,urlcolor=rulered,linkcolor=black}

\begin{document}
\setlength{\parindent}{0pt}
\setlength{\parskip}{4pt}

{\small\color{keywordgray}\textbf{Keywords:}
  time series foundation models \textbullet{}
  zero-shot forecasting \textbullet{}
  benchmarking \textbullet{}
  temporal inductive bias}\par\vspace{4pt}

{\LARGE\bfseries\raggedright
  Foundations without Fundamentals: Zero-Shot Blind Spots in Time Series Foundation Models}\par\vspace{4pt}

{\large\itshape\raggedright
  Chronos-2, Moirai, and Toto Fail Unit Tests for Monotonic Trends, Periodic Signals,
  and Leading Indicators That Classical Models Pass Trivially}\par\vspace{6pt}

{\small\textbf{By} Nafiseh Ghoroghchian, Haipeng Zhang, Shuyi Han,
  Alex Labach \& George Stein (ServiceNow Research; McGill University)}\par\vspace{2pt}

{\small\color{keywordgray}arXiv:2610.02058 \textbullet{} October 2026}
\par\vspace{2pt}\noindent{\color{rulered}\rule{\columnwidth}{0.6pt}}\vspace{6pt}

Despite strong performance on broad benchmarks, prominent Time Series Foundation Models
fail surprisingly basic unit tests---monotonic trends, periodic signals, and leading-indicator
exploitation---that simple classical models handle trivially, revealing a gap between
pre-training scale and fundamental temporal inductive biases that persists on real-world
sensor data.

\begin{mdframed}[backgroundcolor=boxgray,linecolor=rulered,linewidth=1pt,
    innerleftmargin=6pt,innerrightmargin=6pt,
    innertopmargin=6pt,innerbottommargin=6pt]
\textbf{\small AT A GLANCE}\par\vspace{4pt}
\begin{description}[leftmargin=0pt,labelsep=4pt,itemsep=2pt,parsep=0pt,
    font=\small\bfseries]
  \item[Problem:] TSFMs may lack inductive biases for basic temporal patterns despite
    strong broad-benchmark performance.
  \item[Why it matters:] Monotonic trends, periodicity, and causal leading indicators
    are ubiquitous in industrial, scientific, and healthcare time series.
  \item[Method:] SimpleTimeBench---a diagnostic suite of unit tests for temporal
    primitives (trend, periodicity, leading indicators) at multiple difficulty levels.
  \item[Result:] Chronos-2, Moirai, and Toto frequently fail zero-shot unit tests;
    ARIMA/Theta match or exceed them on basic primitives.
  \item[Evidence:] Synthetic controlled series plus real-world building energy meters and
    meteorological stations.
\end{description}
\end{mdframed}

\section{The Big Picture}

Time Series Foundation Models (TSFMs) trained on massive corpora---Chronos-2 on
billions of time points, Moirai on diverse multi-domain data---have achieved
impressive zero-shot performance on established benchmarks such as Monash and ETT.
The natural inference is that these models have internalized general temporal structure,
including basic primitives like monotonic trends, cyclical seasonality, and causal
leading indicators.

This paper challenges that inference by constructing SimpleTimeBench: a diagnostic
benchmark of unit tests for temporal primitives where any well-calibrated forecaster
should produce near-perfect zero-shot results. If a model with billions of training
examples fails these trivial tests, the gap between pre-training scale and genuine
temporal reasoning is larger than aggregate benchmarks suggest.

\section{Why Existing Approaches Fall Short}

\textbf{Conflation problem:} Aggregate benchmark performance combines hundreds of
heterogeneous series; strong performance on diverse data does not guarantee competence
on any individual temporal primitive.

\textbf{No unit tests:} There is no existing benchmark analogous to software unit tests
that verifies whether a model handles basic temporal building blocks correctly.

\textbf{Invisible covariate gaps:} Existing benchmarks rarely test whether models
exploit available exogenous covariates (leading indicators), so systematic failure
to use causal information goes undetected.

\section{Core Mechanism}

\textbf{SimpleTimeBench} comprises three test families:
\begin{itemize}[itemsep=2pt,parsep=0pt]
  \item \emph{Monotonic trend:} Univariate series $x_t = \alpha t + \sigma\epsilon_t$
    with varying slope $\alpha$ and noise $\sigma$. The oracle continues the trend;
    reversals or early flattening constitute failures.
  \item \emph{Periodic signals:} Series with known periods (daily/weekly/annual)
    and controlled amplitudes. A model with temporal periodicity inductive bias
    should closely track cycles.
  \item \emph{Leading indicators:} Multivariate series where channel $x_{t-k}$
    Granger-causes $y_t$ with a known lag $k$, providing complete information
    for an accurate forecast.
\end{itemize}

Each family spans multiple noise levels, slopes, periods, and lag lengths to
characterize sensitivity rather than produce a binary pass/fail verdict.

\section{Technical Formulation}

For the monotonic trend test at noise level $\sigma$ and slope $\alpha$:
\[
x_t = \alpha t + \sigma \epsilon_t, \quad \epsilon_t \sim \mathcal{N}(0,1)
\]
The oracle forecast for horizon $h$ is $\hat{x}_{T+h} = \alpha(T+h)$, achieving MASE $\approx 1$.
A model \emph{passes} if $\text{MASE}_{\text{model}} / \text{MASE}_{\text{oracle}} < \delta$
for threshold $\delta = 2.0$.

For the leading indicator test with lag $k$:
\[
y_t = \phi(x_{t-k}) + \sigma\epsilon_t
\]
The oracle MSE is $\sigma^2$. The leading-indicator utilization rate $\rho$ is:
\[
\rho = 1 - \frac{\text{MSE}_{\text{model}} - \sigma^2}{\text{MSE}_{\text{no-cov}} - \sigma^2}
\]
where $\text{MSE}_{\text{no-cov}}$ is the MSE of a univariate model that ignores $x$.
A $\rho$ near 1 indicates full utilization; TSFMs average $\rho < 0.4$.

\section{Learning or Inference Procedure}

SimpleTimeBench is evaluation-only. The procedure:
\begin{enumerate}[itemsep=2pt,parsep=0pt]
  \item Generate controlled synthetic series at multiple difficulty levels.
  \item Run each TSFM (Chronos-2, Moirai, Toto) in zero-shot mode with
    the model's default context window.
  \item Compute MASE, sMAPE, and $\rho$ for each test family.
  \item Compare against oracle forecasters and classical baselines
    (ARIMA, Theta, ARIMAX).
  \item Validate on real-world sensor datasets containing the tested primitives.
\end{enumerate}

\section{What the Guarantee Says}

SimpleTimeBench provides a falsifiability guarantee: failure on any test family certifies
the existence of natural inputs (with no distributional shift) for which the model
underperforms a trivially simple oracle. This is a constructive, worst-case
characterization---not a claim about average performance.

\section{Experimental Findings}

\begin{itemize}[itemsep=2pt,parsep=0pt]
  \item \textbf{Trend reversal:} Chronos-2 reverses an increasing trend in
    $\approx18\%$ of test cases at low noise; Moirai flattens trends early
    in $\approx23\%$.
  \item \textbf{Periodic signals:} All three models exhibit MASE ratios above 2.0
    relative to the oracle in $>30\%$ of annual-period instances.
  \item \textbf{Leading indicators:} Average utilization rate $\rho < 0.40$
    for all three TSFMs; ARIMAX achieves $\rho > 0.80$ on the same series.
  \item \textbf{Real-world sensors:} On building energy meters, TSFMs consistently
    underperform a tuned Fourier decomposition model on strongly periodic signals.
  \item \textbf{Classical model comparison:} ARIMA and Theta match or exceed TSFMs
    on monotonic trend and leading indicator tests despite being far simpler.
\end{itemize}

\section{Ablations and Interpretation}

\textbf{Context window:} Longer context modestly improves periodic performance
(5--10\% MASE reduction) but does not eliminate failures.

\textbf{Model size:} Larger model variants do not substantially improve
SimpleTimeBench pass rates---the blind spots are not simply a parameter-count issue.

\textbf{Fine-tuning:} A small amount of domain fine-tuning on primitive-containing
examples raises pass rates to $>85\%$, confirming the biases are learnable but
absent in zero-shot pre-trained form.

\section*{Reference}

Nafiseh Ghoroghchian, Haipeng Zhang, Shuyi Han, Alex Labach, George Stein.
\textbf{Foundations without Fundamentals: Zero-Shot Blind Spots in Time Series FMs}.
arXiv:2610.02058, October 2026.\\
\url{https://arxiv.org/abs/2610.02058}

\end{document}
